\subsection{Multiplicity of a module with respect to an ideal}

Let M be a finitely generated A-module and q an ideal of A contained in the radical of A such that M/qM has finite length. Suppose M is not reduced to 0, and set \(d = \dim_{\mathbf{A}}(\mathbf{M})\). By § 4, no. 3, corollary of Th. 2 and no. 4, Remark 2, there exists a unique integer \(e_{\mathfrak{q}}(\mathbf{M}) > 0\) such that, for every integer \(n \geqslant 1\)

\[
\operatorname{long} _ {\mathbf {A}} (\mathrm{M} / \mathfrak {q} ^ {n + 1} \mathrm{M}) = e _ {\mathfrak {q}} (\mathrm{M}) \frac {n ^ {d}}{d !} + \beta_ {n} n ^ {d - 1}
\]

where \(\beta_{n}\) tends to a limit as \(n\) increases indefinitely.

DEFINITION 1. --- The integer \(e_{\mathfrak{q}}(\mathbf{M})\) is called the multiplicity of the A-module \(\mathbf{M}\) relative to the ideal q.

It is also denoted \(e_{\mathfrak{q}}^{\mathbf{A}}(\mathbf{M})\) when one wishes to mention the ring A. When A is local with maximal ideal m, one writes \(e(\mathbf{M})\) or \(e^{\mathbf{A}}(\mathbf{M})\) for \(e_{\mathfrak{m}}(\mathbf{M})\) or \(e_{\mathfrak{m}}^{\mathbf{A}}(\mathbf{M})\) .

Remarks. --- 1) If q' is an ideal of A contained in the radical of A and containing q, then \(e_{\mathfrak{q}'}(\mathbf{M}) \leqslant e_{\mathfrak{q}}(\mathbf{M})\) and, if the \(q'\)-adic filtration of M is q-good, we have \(e_{\mathfrak{q}'}(\mathbf{M}) = e_{\mathfrak{q}}(\mathbf{M})\) (§ 4, no. 3, th. 2).

\begin{enumerate}
\def\labelenumi{\arabic{enumi})}
\setcounter{enumi}{1}
\item
  If M is of finite length, we have \(e_{\mathfrak{q}}(\mathbf{M}) = \mathrm{long}_{\mathbf{A}}(\mathbf{M})\) (§ 4, no. 3, remark 3).
\item
  If \(d > 0\) , one has
\end{enumerate}

\[
\operatorname{long} _ {\mathrm{A} / \mathfrak {q}} (\mathfrak {q} ^ {n} \mathrm{M} / \mathfrak {q} ^ {n + 1} \mathrm{M}) = e _ {\mathfrak {q}} (\mathrm{M}) \frac {n ^ {d - 1}}{(d - 1) !} + \alpha_ {n} n ^ {d - 2}
\]

where \(\alpha_{n}\) tends to a limit as \(n\) increases indefinitely (§ 4, no. 3, corollary to th. 2).

\begin{enumerate}
\def\labelenumi{\arabic{enumi})}
\setcounter{enumi}{3}
\tightlist
\item
  One can reduce the computation of multiplicities to the case where A is local, since, by § 4, no. 4, corollary to th. 3, we have
\end{enumerate}

\[
e _ {\mathrm{q}} (\mathbf {M}) = \sum e _ {\mathrm{q} _ {\mathrm{m}}} (\mathbf {M} _ {\mathrm{m}})
\]

where the summation is extended over the maximal ideals m of A such that

\[
\mathfrak {m} \in \operatorname{Supp} (\mathbf {M}) \cap \mathrm{V} (\mathfrak {q}) \quad \text { and } \quad \dim_ {\mathrm{A} _ {\mathfrak {m}}} (\mathrm{M} _ {\mathfrak {m}}) = d  .
\]

It follows from Remarks 2 and 3 that \(e_{\mathfrak{q}}(\mathbf{M})\) depends only on the graded A/q-module \(\mathrm{gr}_{\mathfrak{q}}(\mathbf{M})\). Consequently:

PROPOSITION 1.--- Let \(\hat{\mathbf{A}}\) and \(\hat{\mathbf{M}}\) be the completions of \(\mathbf{A}\) and \(\mathbf{M}\) for their q-adic topologies; then \(e_{\mathfrak{q}}^{\mathbf{A}}(\mathbf{M}) = e_{\mathfrak{q}\hat{\mathbf{A}}}^{\hat{\mathbf{A}}}(\hat{\mathbf{M}})\) .

PROPOSITION 2. --- Suppose A is regular local (§ 5, no. 1, def. 1); then e(A) = 1.

This follows from th. 1 of § 5, no. 2.

Remark 5. --- One can have \(e(A) = 1\) without A being regular (p.~104, exerc. 5). In fact, a Noetherian local ring A is regular if and only if \(\hat{A}\) is an integral domain and if one has \(e(A) = 1\) (p.~108, exerc. 24).

Example. --- By definition, we have \(e_{\mathfrak{q}^{r}}(\mathbf{M}) = r^{d}e_{\mathfrak{q}}(\mathbf{M})\) where \(d = \dim_{\mathbf{A}}(\mathbf{M})\) . Consequently, if A is regular local, we have \(e_{\mathfrak{m}_{\mathbf{A}}^{r}}(\mathbf{A}) = r^{d}\) . For example, if A is a discrete valuation ring, we have \(e_{\mathfrak{q}}(\mathbf{A}) = \text{long}(\mathbf{A}/\mathfrak{q})\) .

Let q be an ideal of A contained in the radical of A, and let \(\mathcal{C}(\mathfrak{q})\) be the set of classes of finitely generated A-modules M such that \(M/qM\) has finite length. For every \(d \in \mathbf{N}\), denote by \(\mathcal{C}(\mathfrak{q})_{\leq d}\) the part of \(\mathcal{C}(\mathfrak{q})\) consisting of the classes of A-modules of dimension \(\leqslant d\). One defines a map \(e_{\mathfrak{q},d}: \mathcal{C}(\mathfrak{q})_{\leq d} \to \mathbf{Z}\) by \(e_{\mathfrak{q},d}(\mathbf{M}) = e_{\mathfrak{q}}(\mathbf{M})\) if \(\dim(\mathbf{M}) = d\), and \(e_{\mathfrak{q},d}(\mathbf{M}) = 0\) otherwise. This map is additive by prop. 5 of § 4, no. 3; from it one deduces (A, VIII, § 10, no. 2) a homomorphism, still denoted \(e_{\mathfrak{q},d}\), of the Grothendieck group \(K(\mathcal{C}(\mathfrak{q})_{\leq d})\) into \(\mathbf{Z}\), which vanishes on \(K(\mathcal{C}(\mathfrak{q})_{\leq d-1})\). Reasoning as in § 1, no. 5, we deduce:

PROPOSITION 3. — Let \(\mathbf{M}\) be a finitely generated A-module, of dimension \(d \geqslant 0\). Let \(\Phi\) be the set of minimal elements \(\mathfrak{p}\) of \(\operatorname{Supp}(\mathbf{M})\) such that \(\dim(\mathbf{A}/\mathfrak{p}) = d\). Let \(\mathfrak{q}\) be an ideal of A, contained in the radical of A, and such that \(\mathbf{M}/\mathfrak{q}\mathbf{M}\) is of finite length. We have

\[
e _ {\mathfrak {q}} (\mathrm{M}) = \sum_ {\mathfrak {p} \in \Phi} \operatorname{long} _ {\mathrm{A} _ {\mathfrak {p}}} (\mathrm{M} _ {\mathfrak {p}}) \cdot e _ {\mathfrak {q}} (\mathrm{A} / \mathfrak {p}).
\]

COROLLARY. — Suppose A is semi-local and let q be an ideal of definition of A.

\begin{enumerate}
\def\labelenumi{\alph{enumi})}
\item
  We have \(e_{\mathfrak{q}}(\mathbf{A}) = \sum_{\mathfrak{p}} e_{\mathfrak{q}}(\mathbf{A}/\mathfrak{p})\) , where \(\mathfrak{p}\) ranges over the set of minimal prime ideals of \(\mathbf{A}\) such that \(\dim(\mathbf{A}/\mathfrak{p}) = \dim(\mathbf{A})\) .
\item
  Suppose A is integral and let M be a finitely generated A-module such that \(\dim_{\mathbf{A}}(\mathbf{M}) = \dim(\mathbf{A})\) . Then one has \(e_{\mathfrak{q}}(\mathbf{M}) = \operatorname{rg}(\mathbf{M}) \cdot e_{\mathfrak{q}}(\mathbf{A})\) .
\end{enumerate}

